\documentclass[11pt, a4paper]{article}

% --- 基础宏包配置 ---
\usepackage[UTF8]{ctex}              % 中文支持
\usepackage[margin=2.5cm]{geometry}  % 页面边距
\usepackage{amsmath, amssymb, amsfonts} % 数学公式
\usepackage{graphicx}                % 图片支持
\usepackage{xcolor}                  % 颜色支持
\usepackage{hyperref}                % 超链接
\usepackage{enumitem}                % 列表定制
\usepackage{booktabs}                %以此画出更专业的表格线
\usepackage{tabularx}                % 自动调整宽度的表格
\usepackage{listings}                % 代码环境
\usepackage{tcolorbox}               % 文本框（用于案例展示）
\usepackage{fancyhdr}                % 页眉页脚

% --- 样式设置 ---
\hypersetup{
    colorlinks=true,
    linkcolor=blue,
    filecolor=magenta,      
    urlcolor=cyan,
}

% 定义代码块样式 (仿 Markdown 风格)
\definecolor{codegray}{rgb}{0.95,0.95,0.95}
\definecolor{keywordcolor}{rgb}{0.0, 0.0, 0.5}
\lstdefinestyle{mystyle}{
    backgroundcolor=\color{codegray},   
    basicstyle=\ttfamily\scriptsize,
    breakatwhitespace=false,         
    breaklines=true,                 
    captionpos=b,                    
    keepspaces=true,                 
    showspaces=false,                
    showstringspaces=false,
    showtabs=false,                  
    tabsize=2,
    frame=none,
    keywordstyle=\color{keywordcolor}\bfseries,
    morekeywords={算法, For, While, If, Else, Return, Break, in, range, len, list, append, add, True, False, None},
    commentstyle=\color{gray}
}
\lstset{style=mystyle}

% 紧凑列表设置
\setlist{nosep}

% 标题设置
\title{\textbf{Tokenizer 深度解析与算法汇总}}
\author{}
\date{}

\begin{document}

\maketitle

\section{Tokenizer 定义}

Tokenizer 是字符串到整数数组的双射 $f$。例如我们想要一个 "Tokenizer is amazing!" 到整数数组的对应，则需要设计一个算法，支持将其转换成一个整数数组，并由 $f^{-1}$ 可以将整数数组再还原回字符串。$f$ 即是 encoder，$f^{-1}$ 即是 decoder。Tokenizer 是人类和 LLM 中的“翻译家”。

\begin{tcolorbox}[colback=blue!5!white, colframe=blue!75!black, title=分词方案示例]
\textbf{方案一：}
\begin{enumerate}
    \item 统计整个字符串中有哪些不同的字符，['o', 'z', 'm', 'n', ' ', 'a', 'i', '!', 'r', 'g', 'k', 's', 'e', 'T']；
    \item 构建字符和索引的关系映射表，并且每个字符用索引号来表示；\\
    关系表：\{'o':0, 'z':1, 'm':2, 'n':3, ' ':4, 'a':5, 'i':6, '!':7, 'r':8, 'g':9, 'k':10, 's':11, 'e':12, 'T':13\} \\
    整数数组：[13, 0, 10, 12, 3, 6, 1, 12, 8, 4, 6, 11, 4, 5, 2, 5, 1, 6, 3, 9, 7]
    \item 整数数组反查关系表还原字符串；
\end{enumerate}

\textbf{方案二：}
\begin{enumerate}
    \item 按照空格和标点划分字符，['Tokenizer', 'is', 'amazing', '!', ' ']；
    \item 构建字符和索引的关系映射表，并且每个字符用索引号来表示；\\
    关系表：\{'Tokenizer':0, 'is':1, 'amazing':2, '!':3, ' ':4\} \\
    整数数组：[0, 4, 1, 4, 2, 3]
    \item 整数数组反查关系表还原字符串；
\end{enumerate}
\end{tcolorbox}

这几个方案的区别就是按什么粒度去拆分字符，拆分后的字符/字符串给上唯一数字。如果我们现在想预先构建一个关系表，对于所有传入的字符都能按照关系表进行拆分，拆分后转化为数字，那么这个关系表该怎么构建？

OpenAI 的 tokenizer 可视化平台：\url{https://platform.openai.com/tokenizer}

我们定义一个长度为 $k$ 的整数数组 $\in \Sigma^{k}$，其中 $\Sigma\subset \mathbb{Z}$，$\# \Sigma$ 是我们通常所说的 Tokens 数量。

\section{Tokenizer 重要性}

\begin{itemize}
    \item \textbf{影响模型性能}：好的 Tokenizer 设计可以提升模型的语义理解能力
    \item \textbf{决定上下文长度}：模型的最大上下文窗口是按 token 数量计算的
    \item \textbf{关系推理效率}：token 数量直接影响模型的推理速度和计算成本
    \item \textbf{支持多语言}：优秀的 Tokenizer 设计可以更好地处理不同语言的特性
\end{itemize}

但是 Tokenizer 的选择在很大程度上是“玄学”和“炼丹”。没人能拿出一套数学证明说：“对于语言 $L$ 和任务 $T$，做法 A（如 WordPiece）的误差上界一定比做法 B（BPE）低。”

\section{Tokenizer 核心内容与分类}

\begin{enumerate}
    \item 如何离线构建 token 词表（构建）；
    \item 有词表后，如何对输入的字符串进行快速拆分，并转为 token id（检索）；
\end{enumerate}

好的分词应该是能够恰当粒度地拆分 token。

\subsection*{字符级分词}
前面方案一采用了字符级进行分词，就是以 char 为最小粒度。
\begin{itemize}
    \item \textbf{优点}：词表极小，比如：26 个英文字母几乎可以组合出所有词，5000 多个中文常用字基本也能组合出足够的词汇。
    \item \textbf{缺点}：无法承载丰富的语义，英文中尤为明显，但中文却是较为合理，中文中用此种方式较多；序列长度大幅增长。
\end{itemize}

\subsection*{词组级分词}
算法题中方案二采用了词组级进行分词，就是以一个词、空格、标点为最小粒度。
\begin{itemize}
    \item \textbf{优点}：词的边界和含义得到保留。
    \item \textbf{缺点}：词表大，稀有词学不好；存在 OOV 问题（Out-of-Vocabulary）。
\end{itemize}

\subsection*{子词级分词}
为了平衡以上两种方法，又提出了基于 subword 进行分词：它可以较好的平衡词表大小与语义表达能力，子词级分词算法在现代大模型中应用最为广泛。常见的子词算法有 Byte-Pair Encoding (BPE)、WordPiece、SentencePiece 等。

\section{BPE（Byte Pair Encoding）算法}

BPE 是一种基于统计的子词分词算法，由 Philip Gage 于 1994 年提出，最初用于数据压缩。2016 年被引入 NLP 领域，成为现代大模型的标准分词方法。核心思想是不断合并“最常一起出现的相邻符号对”，逐步形成子词。

\begin{lstlisting}[caption={构建 BPE 算法}]
算法 构建BPE(语料库 text_corpus, 目标词表大小 target_vocab_size):
    词频表 = 统计单词频率(text_corpus)
    词表 = 获取所有基础字符(词频表) 
    For 单词 in 词频表:
        # 将 "hug" 映射为 ["h", "u", "g"]
        当前词拆分 = 将单词拆分为字符列表(单词)
        词拆分字典[单词] = 当前词拆分
        合并规则 = []  # 用于存储 ('a', 'bc') -> 'abc' 这样的规则
        While (len(词表) < target_vocab_size):
        所有字符对 = {}
            For 单词, 频率 in 词频表:
            当前拆分 = 词拆分字典[单词]
                For i in range(len(当前拆分) - 1):
                    对 = (当前拆分[i], 当前拆分[i+1])
                    所有字符对[对] += 频率
        If 所有字符对 为空: Break
        最佳对 = KeyWithMaxVal(所有字符对) # 任选一对频率最高的
        合并规则.add(最佳对), 词表.add(最佳对)
        For 单词 in 词拆分字典:
            词拆分字典[单词] = MergePair(词拆分字典[单词], 最佳对)
    Return 合并规则
\end{lstlisting}

\begin{lstlisting}[caption={BPE 编码算法}]
算法 BPE编码(输入文本 input_text, 合并规则 merge_rules, 词表 vocab):
    TokenID列表 = []
    原始单词列表 = 分词(input_text)
    For 单词 in 原始单词列表:
        当前拆分 = 将单词拆分为字符列表(单词)
        While (len(当前拆分) > 1):
            最佳对 = None
            最小索引 = 无穷大
            For i in range(len(当前拆分) - 1):
                对 = (当前拆分[i], 当前拆分[i+1])
                If 对 in 合并规则:
                    当前索引 = IndexOf(合并规则, 对)
                    If 当前索引 < 最小索引:
                        最佳对 = 对
                        最小索引 = 当前索引
            If 最佳对 为空: Break
            当前拆分 = MergePair(当前拆分, 最佳对)
        For 片段 in 当前拆分:
            TokenID列表.append(vocab[片段] if 片段 in vocab else vocab["<UNK>"])
    Return TokenID列表
\end{lstlisting}

\section{WordPiece 算法}

WordPiece 是 Google 开发的子词分词算法，在 BPE 的基础上进行了改进，主要用于 BERT 系列模型。它是基于概率的子词构建算法，合并逻辑从简单的频率最大化（BPE）转为选择点互信息（PMI）最大的子词对，从而最大化训练数据的对数似然（假设 word 独立同分布）。

\begin{lstlisting}[caption={构建 WordPiece}]
算法 构建WordPiece(语料库 text_corpus, 目标词表大小 target_vocab_size):
    词频表 = 统计单词频率(text_corpus)
    词表 = 获取所有基础字符(词频表) 
    For 单词 in 词频表:
        # 将 "hug" 映射为 ["h", "##u", "##g"]
        当前词拆分 = 将单词拆分为字符列表(单词)
        词拆分字典[单词] = 当前词拆分
    While (len(词表) < target_vocab_size):
        符号计数 = DefaultDict(0)
        对计数 = DefaultDict(0)
        For 单词, 频率 in 词频表:
            当前拆分 = 词拆分字典[单词]
            For 符号 in 当前拆分:
                符号计数[符号] += 频率
            For i in range(len(当前拆分) - 1):
                对 = (当前拆分[i], 当前拆分[i+1])
                对计数[对] += 频率
        最佳对 = None
        最高分数 = -1
        For 对 in 对计数:
            left, right = 对
            分数 = 对计数[对] / (符号计数[left] * 符号计数[right])
            If 分数 > 最高分数:
                最高分数 = 分数
                最佳对 = 对
        If 最佳对 is None: Break
        词表.add(最佳对)
        For 单词 in 词拆分字典:
            词拆分字典[单词] = MergePair(词拆分字典[单词], 最佳对)
    Return 词表
\end{lstlisting}

\begin{lstlisting}[caption={WordPiece 编码}]
算法 WordPiece编码(输入文本 input_text, 词表 vocab):
    TokenID列表 = []
    原始单词列表 = 基础分词(input_text) # 如按空格和标点切分
    For 单词 in 原始单词列表:
        # 确保当前词没有过长，否则 <UNK>
        字符列表 = list(单词)
        当前单词Tokens = []
        起始位置 = 0
        是否为未知词 = False
        While 起始位置 < len(字符列表):
            结束位置 = len(字符列表)
            匹配到的Token = None
            While 结束位置 > 起始位置:
                子串 = 单词[起始位置 : 结束位置]
                if 起始位置 > 0:
                    子串 = "##" + 子串
                if 子串 in vocab:
                    匹配到的Token = 子串
                    Break
                结束位置 -= 1
            If 匹配到的Token is None:
                是否为未知词 = True
                Break
            当前单词Tokens.append(匹配到的Token)
            起始位置 = 结束位置
        If 是否为未知词:
            TokenID列表.append(vocab["<UNK>"])
        Else:
            For token in 当前单词Tokens:
                TokenID列表.append(vocab[token])
    Return TokenID列表
\end{lstlisting}

\section{Unigram Language Model (ULM) 算法}

Unigram Language Model (ULM) 是一种基于概率的子词分词算法，由 Taku Kudo 在 2018 年随 SentencePiece 库提出。与 BPE 和 WordPiece 的“自底向上”合并策略不同，ULM 采用“自顶向下”的剪枝策略。它初始化一个巨大的词表，然后通过 EM 算法迭代优化，逐步剔除对总似然值（Log-Likelihood）贡献最小的子词，直到达到目标词表大小。它通过 Forward-Backward 算法寻找产生输入句子的最大概率切分路径。

\begin{lstlisting}[caption={构建 ULM}]
算法 构建ULM(语料库 text_corpus, 目标词表大小 target_vocab_size):
    候选词表 = 获取所有子串及字符(text_corpus)
    词概率 = {词: 频率/总频次 For 词, 频率 in 统计频率(候选词表)}
    While (len(候选词表) > target_vocab_size):
        # 工程上只跑一轮 EM 作为近似（“热启动”效应）
        # --- E-Step: 固定词表和概率，统计语料中各子词的预期出现次数 ---
        预期计数 = DefaultDict(0)
        For 句子 in text_corpus:
            # 1. Forward Pass (Alpha): 计算从句首到位置 i 的所有路径概率和
            Alpha[0] = 1
            For t in range(len(句子)):
                For 词 in 匹配字典(句子, start=t):
                    Alpha[t + len(词)] += Alpha[t] * 词概率[词]
            # 2. Backward Pass (Beta): 计算从位置 i 到句尾的所有路径概率和
            Beta[len(句子)] = 1
            For t in range(len(句子)-1, -1, -1):
                For 词 in 匹配字典(句子, start=t):
                    Beta[t] += 词概率[词] * Beta[t + len(词)]
            句子总概率 = Alpha[len(句子)]
            For t in range(len(句子)):
                For 词 in 匹配字典(句子, start=t):
                    后验 = (Alpha[t] * 词概率[词] * Beta[t + len(词)]) / 句子总概率
                    预期计数[词] += 后验
        # M-Step: 更新概率：每个词的后验概率归一化分配为每个词的概率
        总预期频次 = Sum(预期计数.values())
        For 词, 频次 in 预期计数.items():
            词概率[词] = 频次 / 总预期频次
        # --- 剪枝：移除对总似然贡献最小的词 ---
        损失列表 = []
        当前总似然 = 计算语料库总似然(text_corpus, 词概率)
        For 词 in 候选词表:
            # 估算移除该词后，总似然会下降多少（Loss）
            损失 = 计算移除损失(词, 预期计数) #Viterbi，实现见ULM编码
            损失列表.append((词, 损失))
        排序(损失列表, key=损失)
        移除数量 = max(len(候选词表) * 0.2, len(候选词表) - target_vocab_size)
        要移除的词 = 损失列表[:移除数量]
        For 词, _ in 要移除的词:
            If 词 not in 基础字符集:
                候选词表.remove(词)
                Del 词概率[词]
    Return 候选词表, 词概率
\end{lstlisting}

\begin{lstlisting}[caption={ULM 编码 (Viterbi)}]
算法 ULM编码(输入文本 input_text, 词表 vocab, 词概率 probabilities):
    N = len(input_text)
    DP = [-无穷大] * (N + 1)
    DP[0] = 0.0
    回溯指针 = [0] * (N + 1)
    For i in range(1, N + 1):
    For j in range(0, i):
        子串 = input_text[j:i]
        If 子串 in vocab:
            当前概率 = DP[j] + log(probabilities[子串])
            If 当前概率 > DP[i]:
                DP[i] = 当前概率
                回溯指针[i] = j
     TokenID列表 = []
     当前位置 = N
     While 当前位置 > 0:
         上一位置 = 回溯指针[当前位置]
         片段 = input_text[上一位置:当前位置]
         TokenID列表.insert(0, vocab[片段] if 片段 in vocab else vocab["<UNK>"])
         当前位置 = 上一位置
    Return TokenID列表
\end{lstlisting}

\section{BPE、WordPiece、Unigram 对比}

\begin{table}[h]
\centering
\small
\renewcommand{\arraystretch}{1.5}
\begin{tabularx}{\textwidth}{|l|X|X|p{2.5cm}|X|}
\hline
\textbf{算法} & \textbf{核心机制} & \textbf{主要解决问题} & \textbf{应用场景} & \textbf{关键差异} \\ \hline
\textbf{BPE} & 基于\textbf{频率统计}合并高频字符对，逐步构建词表 & 减少词表大小，覆盖低频词和未登录词（OOV），避免因词表过大或过小导致的模型效率问题 & 机器翻译、RoBERTa 等模型 & 合并策略基于字符对频率，可能导致歧义切分；需预先设定词表大小 \\ \hline
\textbf{WordPiece} & 通过计算合并后对语言模型\textbf{似然值的提升}选择字符对，采用贪心最长匹配策略 & 优化子词组合的语义独立性，提升模型对上下文的理解能力 & BERT、DistilBERT 等预训练模型（没那么关心 OOV） & 合并策略基于似然增益，更注重语义关联性；与 BPE 相比生成更稳定的子词组合 \\ \hline
\textbf{Unigram} & 基于概率模型\textbf{反向剪枝}，从大词表开始\textbf{逐步删除低概率子词}，保留高概率子词构建词表 & 平衡子词粒度的灵活性与词表效率，避免局部最优问题 & XLNet、跨语言模型（如 SentencePiece） & 采用概率模型而非合并策略，支持多分词路径选择；无需预分词 \\ \hline
\end{tabularx}
\end{table}

\section{SentencePiece 算法}

SentencePiece 是谷歌推出的子词开源工具包，其中集成了 BPE、ULM 子词算法。除此之外，SentencePiece 还能支持字符和词级别的分词。更进一步，为了能够处理多语言问题，SentencePiece 将句子视为 Unicode 编码序列，从而子词算法不用依赖于语言的表示，支持可逆的编码和解码的过程。

\subsection*{算法对比分析}

\begin{table}[h]
\centering
\begin{tabularx}{\textwidth}{|l|l|l|X|X|l|}
\hline
\textbf{算法} & \textbf{核心思想} & \textbf{代表模型} & \textbf{优点} & \textbf{缺点} & \textbf{适用场景} \\ \hline
BPE & 合并高频字符对 & GPT-2/3/4 & 简单高效，词表可控 & 贪心策略非全局最优 & 生成任务，多语言 \\ \hline
WordPiece & 互信息驱动合并 & BERT系列 & 语义更连贯 & 中文处理需预分词 & 理解任务，短文本 \\ \hline
SentencePiece & 语言无关端到端 & T5, mBART & 支持无空格语言 & 训练复杂 & 多语言 NLP 任务 \\ \hline
\end{tabularx}
\end{table}

\section{BBPE（Byte-level Byte Pair Encoding）}

BBPE 是一种以“字节”作为初始最小单位（e.g. 中->E4 B8 AD）的 BPE 算法，通过不断合并高频相邻字节对，构建子词词表，从而实现对任意文本的无 OOV 分词。初始词表固定 256 byte。

\section{智能分词策略}

BERT + CRF 不只是拆分实体/词性，核心是 “语义+上下文驱动的动态分词”。

\subsection*{可扩展词表}
可扩展词表用于应对日益更新的新词。

\textbf{词表设计原理}
\begin{enumerate}
    \item 每个层级有独立的 ID 空间：
    \begin{itemize}
        \item Level 0 (基础字符): 0-255
        \item Level 1 (常用子词): 256-65535
        \item Level 2 (通用词汇): 65536-1097727
        \item Level 3 (领域词汇): 1097728-33554431
        \item Level 4 (专业术语): 33554432-201326591
    \end{itemize}
    \item 新词只会添加到 Level 3 或 Level 4，底层词汇永久固定
\end{enumerate}

\textbf{新词入库流程（以 Level3~4 为例）}
\begin{quote}
新词检测 $\rightarrow$ 自动分层分类 $\rightarrow$ 智能审核 $\rightarrow$ 入库映射 $\rightarrow$ 效果验证
\end{quote}

\begin{enumerate}
    \item \textbf{新词检测（自动）}
    \begin{itemize}
        \item 触发条件：推理时 tokenizer 遇到未命中词（触发 byte-level 回退），或批量文本分析中高频出现的 N-gram（2-4 元组）；
        \item 算法示例：用「频次阈值 + 领域相关性得分」筛选，比如金融文本中 “数字人民币” 出现频次 $>500$ 次，领域相关性 $>0.8$ 则标记为候选新词。
    \end{itemize}
    
    \item \textbf{自动分层分类（自动）}
    分层规则：基于预设的层级判定模型（轻量 BERT 微调），输入新词 + 上下文，输出层级概率：
    \begin{itemize}
        \item Level3（领域层）：领域内通用（如医疗的 “PD-1 抑制剂”、金融的 “量化对冲”）；
        \item Level4（专业层）：细分场景术语（如 AI 的 “LoRA 适配器”、法律的 “表见代理”）；
    \end{itemize}
    示例：“量子计算芯片” 经模型判定，领域相关性 0.92、通用性 0.3，自动分配到 Level4。

    \item \textbf{智能审核（人工可介入）}
    \begin{itemize}
        \item 自动审核：通过 “语义一致性校验”（新词与层级词汇语义相似度 $>0.7$）、“无冲突校验”（不与已有词重复 / 歧义）；
        \item 人工介入：仅针对高风险场景（如多义新词 “算力租赁”，可归属 AI 或金融领域），人工标注最终层级，无需逐词审核。
    \end{itemize}

    \item \textbf{入库映射（全自动）}
    \begin{itemize}
        \item 分配 ID：按层级 ID 空间规则，在对应层级末尾分配唯一 ID（如 Level4 新增 “量子计算芯片”，ID=33554436）；
        \item 嵌入初始化：用 “领域词向量迁移” 生成初始嵌入（如基于 “芯片”“量子计算” 的向量加权求和），避免随机初始化。
    \end{itemize}
    
    \item \textbf{模型轻量微调}
    词表新增 token 后，可以在小数据集上，锁层或者 LoRA 微调模型来适配新 token。
\end{enumerate}

\section{多模态 Tokenizer}

将文本、图像、音频、视频等不同模态的数据统一映射为“同一离散 token 空间”，使 Transformer 能在不区分模态的情况下进行建模与对齐。（毕竟 Transformer 模型是针对离散数据发明的）

图像模态处理：从 Patch 到 Token，连续表示离散化（Vector Quantization, VQ）。

\textbf{VQ 算法：}
\begin{itemize}
    \item 建立一个 codebook（视觉的词表，学习方法 K-means、VQ-VAE）
    \item 每个 patch 向量到最近的 codebook 索引（e.g. patch\_1 $\rightarrow$ token\_1024, patch\_2 $\rightarrow$ token\_587）， 这些索引就是“图像 token”。
\end{itemize}

\section{DeepSeekOCR}

DeepSeekOCR“靠视觉压缩一切”，通过用少量的视觉 token 来表示原本需要大量文本 token 的内容，以此降低大模型的计算开销，解决了大模型处理长文本的算力爆炸难题，通过视觉压缩实现了 10-20 倍的效率提升。

视觉 Tokenizer 给文本 Tokenizer 的启示：目前的文本 token 并不能充分地利用 LLM 模型内部给每个 token 预留的表示能力。

\section{各 LLM 使用的 Tokenizer}

\begin{table}[h]
\centering
\begin{tabularx}{\textwidth}{|l|l|l|X|}
\hline
\textbf{模型系列} & \textbf{Tokenizer类型} & \textbf{词汇表大小} & \textbf{特色功能} \\ \hline
GPT系列 & BBPE & 50,257 & 支持任意 UTF-8 字符，字节级处理 \\ \hline
Qwen系列 & BBPE & 150,000+ & 中文优化，大词表设计 \\ \hline
DeepSeek & BBPE & 128,000 & 多语言支持，代码理解优化 \\ \hline
Doubao & BBPE + 视觉编码 & 150,000+ & 多模态统一编码 \\ \hline
Gemini & 多模态统一Tokenizer & 未知 & 图文音视频统一处理 \\ \hline
Claude & 自定义Tokenizer & 未知 & 长文本处理优化 \\ \hline
\end{tabularx}
\end{table}

\end{document}